This appendix gives the proofs supporting the geometric interpretation and structural analysis in the main paper. It establishes a directional relation for the GST projection before normalization and bias, the concentration envelope of tangent-metric attention, and the factorized proximal structure of the temporal affinity.
\section{Blockwise Directional Equivalence on a Product of Spheres}
\label{app:blockwise}
\textbf{Proposition S.1:} Let \(K\geq1\) and \(d\geq2\). For every domain \(k\), assume \(\|\overline{\mathbf p}_k\|_2\geq\epsilon\) and let the gated vector in~\eqref{eq:domain-alignment} satisfy \(\mathbf h_t^{(k)}\neq\mathbf0\), and define
\label{prop:block-direction}
\begin{align}
\boldsymbol{\upsilon}_t^{(k)}
&=\frac{\mathbf h_t^{(k)}}{\|\mathbf h_t^{(k)}\|_2},
&\theta_{t,k}
&=\arccos\!\left((\boldsymbol{\upsilon}_t^{(k)})^\top\mathbf p_k\right),\nonumber\\
\mathbf u_t^{(k)}
&=\Pi_{\mathbf p_k}\mathbf h_t^{(k)},
&\omega_{t,k}
&=\frac{\theta_{t,k}}
{\|\mathbf h_t^{(k)}\|_2\sin\theta_{t,k}}.
\label{eq:app-block-definitions}
\end{align}
Assume \(\theta_{t,k}\in(0,\pi)\) for all \(k\). On the product manifold
\(\mathcal P=(\Sph^{d-1})^K\), stack the base points, normalized aligned vectors, and tangent components as \(\mathbf p=(\mathbf p_1,\ldots,\mathbf p_K)\), \(\boldsymbol{\upsilon}_t=(\boldsymbol{\upsilon}_t^{(1)},\ldots,\boldsymbol{\upsilon}_t^{(K)})\), and \(\mathbf u_t=[(\mathbf u_t^{(1)})^\top,\ldots,(\mathbf u_t^{(K)})^\top]^\top\in\R^{Kd}\). Then
\begin{equation}
\operatorname{Log}_{\mathbf p}(\boldsymbol{\upsilon}_t)
=\boldsymbol{\Omega}_t\mathbf u_t,
\qquad
\boldsymbol{\Omega}_t
=\operatorname{blkdiag}_{k=1}^{K}(\omega_{t,k}\mathbf I_d)\succ\mathbf0.
\label{eq:app-log}
\end{equation}
Hence each pre-normalization projected block is collinear with, and points in the same direction as, the corresponding spherical logarithmic map. The result characterizes the projected direction. With the implemented bias, \(\mathbf z_t^{(k)}-\mathbf b_k\in\Tangent_{\mathbf p_k}\Sph^{d-1}\), so the final token lies in the corresponding affine tangent space.

\emph{Proof:} The product-manifold logarithmic map is the concatenation of its component maps. For domain \(k\), the logarithmic map on \(\Sph^{d-1}\) is
\begin{equation}
\log_{\mathbf p_k}(\boldsymbol{\upsilon}_t^{(k)})
=\frac{\theta_{t,k}}{\sin\theta_{t,k}}
\left(\boldsymbol{\upsilon}_t^{(k)}
-\cos\theta_{t,k}\,\mathbf p_k\right).
\label{eq:app-log-map}
\end{equation}
Because
\(\cos\theta_{t,k}=(\mathbf h_t^{(k)})^\top\mathbf p_k/
\|\mathbf h_t^{(k)}\|_2\),
\begin{equation}
\boldsymbol{\upsilon}_t^{(k)}
-\cos\theta_{t,k}\,\mathbf p_k
=\frac{\Pi_{\mathbf p_k}\mathbf h_t^{(k)}}
{\|\mathbf h_t^{(k)}\|_2}
=\frac{\mathbf u_t^{(k)}}{\|\mathbf h_t^{(k)}\|_2}.
\label{eq:app-tangent}
\end{equation}
Substitution into~\eqref{eq:app-log-map} gives
\(\log_{\mathbf p_k}(\boldsymbol{\upsilon}_t^{(k)})
=\omega_{t,k}\mathbf u_t^{(k)}\). Stacking the \(K\) identities yields~\eqref{eq:app-log}. Since every \(\omega_{t,k}>0\), the direction of each nonzero block is preserved. At \(\theta_{t,k}=0\), both component vectors vanish and no direction is defined; at \(\theta_{t,k}=\pi\), the component logarithmic map is not unique.

The directional agreement follows because both constructions use the same orthogonal component \(\Pi_{\mathbf p_k}\mathbf h_t^{(k)}\), and spherical curvature multiplies it by a positive blockwise scale. Normalization preserves the direction of each nonzero component and controls its norm by \(\sqrt d\). The learned bias then translates the tangent representation, as specified in~\eqref{eq:token}.

\section{Concentration Analysis of Tangent-Metric Attention}
\label{app:attention-proof}
\textbf{Proposition S.2:} For head \(h\), let \(\|\mathbf p_h\|_2\leq1\) as ensured by the implemented normalization, and define
\label{prop:attention-concentration}
\begin{equation}
D_{\max}^{h}
=\max_{1\leq i,j\leq K}\|\mathbf q_i^h-\mathbf k_j^h\|_2^2
\leq4.
\label{eq:app-diameter}
\end{equation}
Then the logits and weights in~\eqref{eq:attention-logit}--\eqref{eq:attention-weight} satisfy
\begin{align}
-\gamma_hD_{\max}^{h}
&\leq \ell_{ij}^{h}\leq0,
\label{eq:app-logit-bound}\\
\max_j a_{ij}^{h}
&\leq\frac{1}{1+(K-1)e^{-\gamma_hD_{\max}^{h}}},
\label{eq:app-weight-bound}\\
\mathsf H(\mathbf a_{i,\cdot}^{h})
&\geq\log\!\left(
1+(K-1)e^{-\gamma_hD_{\max}^{h}}
\right).
\label{eq:app-entropy-bound}
\end{align}

\emph{Proof:} By~\eqref{eq:qkv},
\(\|\bar{\mathbf q}_i^h\|_2\leq1\) and
\(\|\bar{\mathbf k}_j^h\|_2\leq1\). The operator \(\mathbf I-\mathbf p_h\mathbf p_h^\top\) has spectral norm at most one for \(\|\mathbf p_h\|_2\leq1\), so
\(\|\mathbf q_i^h\|_2\leq1\) and
\(\|\mathbf k_j^h\|_2\leq1\). The triangle inequality gives
\(\|\mathbf q_i^h-\mathbf k_j^h\|_2\leq2\), proving~\eqref{eq:app-diameter};~\eqref{eq:app-logit-bound} then follows from \(\gamma_h>0\).

Let \(j^\star=\argmax_j \ell_{ij}^{h}\). The softmax can be written as
\begin{equation}
\max_j a_{ij}^{h}
=\frac{1}{1+
\sum_{m\neq j^\star}
e^{\ell_{im}^{h}-\ell_{ij^\star}^{h}}}.
\label{eq:app-softmax}
\end{equation}
The logit interval implies
\(\ell_{im}^{h}-\ell_{ij^\star}^{h}\geq-\gamma_hD_{\max}^{h}\). Lower-bounding every term in the denominator proves~\eqref{eq:app-weight-bound}. Finally, for any discrete distribution \(\mathbf a\),
\(\mathsf H(\mathbf a)\geq-\log\max_j a_j\), which gives~\eqref{eq:app-entropy-bound}.

For finite \(\gamma_h\) and \(K>1\), the peak-weight bound is strictly below one and the entropy bound is positive. The proposition therefore gives an explicit concentration envelope for tangent-metric attention at a finite inverse bandwidth. On the unit-normalization branch, the head-specific operator is a rank-\((d_h-1)\) orthogonal projector, selecting complementary tangent directions for cross-domain relations.

\section{Raw-Affinity Rank Budget, Proximal Sparsity, and Residual-Output Bound}
\label{app:affinity}
Let \(\mathbf F\in\R^{T\times d_f}\),
\(\mathbf W_u,\mathbf W_v\in\R^{r_\ell\times d_f}\), and define
\begin{alignat}{2}
\mathbf U
&=\mathbf F\mathbf W_u^\top
&\qquad&\in\R^{T\times r_\ell},
\label{eq:app-u}\\
\mathbf V
&=\mathbf F\mathbf W_v^\top
&&\in\R^{T\times r_\ell},
\label{eq:app-v}\\
\mathbf C_{\mathrm{raw}}
&=\frac{1}{\sqrt{r_\ell}}\mathbf U\mathbf V^\top
&&\in\R^{T\times T},
\label{eq:app-craw}\\
\mathbf C
&=\ReLU\!\left(
\mathbf C_{\mathrm{raw}}-\vartheta\mathbf1_T\mathbf1_T^\top
\right)
&&\in\R^{T\times T}.
\label{eq:app-c}
\end{alignat}

\textbf{Proposition S.3.1:} The pre-threshold affinity satisfies
\begin{equation}
\rank(\mathbf C_{\mathrm{raw}})\leq r_\ell,
\label{eq:app-rank}
\end{equation}
independent of \(T\).

\emph{Proof:} Since \(\mathbf U,\mathbf V\in\R^{T\times r_\ell}\),
\begin{equation*}
\rank(\mathbf U\mathbf V^\top)
\leq\min\{\rank(\mathbf U),\rank(\mathbf V)\}
\leq r_\ell.
\end{equation*}
Multiplication by \(1/\sqrt{r_\ell}\) preserves rank, and the bound is attained when both factors have full column rank. The elementwise threshold in~\eqref{eq:app-c} can increase matrix rank, so Proposition~S.3.1 applies only to the raw affinity. The proposition characterizes the compact pre-threshold relation, while Propositions~S.3.2--S.3.3 characterize the nonnegative sparse support and residual aggregation used by the network. Asymmetric entries encode directional temporal dependence, and diagonal entries retain frame self-information.

\textbf{Proposition S.3.2:} Let
\(g_{\vartheta}(\mathbf G)=\vartheta\|\mathbf G\|_1+\iota_{\geq0}(\mathbf G)\), where \(\|\cdot\|_1\) sums entrywise absolute values and \(\iota_{\geq0}\) is zero on the nonnegative orthant and \(+\infty\) elsewhere. Then
\begin{align}
\mathbf C
&=\argmin_{\mathbf G\geq0}
\frac{1}{2}\|\mathbf G-\mathbf C_{\mathrm{raw}}\|_F^2
+\vartheta\|\mathbf G\|_1,
\label{eq:app-prox-problem}\\
\mathbf C&=\prox_{g_{\vartheta}}(\mathbf C_{\mathrm{raw}}),
\label{eq:app-prox}
\end{align}
at fixed \(\mathbf C_{\mathrm{raw}}\), with
\begin{equation}
\vartheta_2>\vartheta_1
\ \Longrightarrow\
\operatorname{supp}(\mathbf C_{\vartheta_2})
\subseteq\operatorname{supp}(\mathbf C_{\vartheta_1}).
\label{eq:app-support}
\end{equation}

\emph{Proof:} The objective is separable. For one entry \(x=C_{\mathrm{raw},ts}\), solve
\begin{equation*}
\min_{g\geq0}\ \frac12(g-x)^2+\vartheta g.
\end{equation*}
Strict convexity gives the unique solution
\(g^\star=\max\{x-\vartheta,0\}\). Applying this result entry-wise proves~\eqref{eq:app-prox-problem}--\eqref{eq:app-prox}. Since \(C_{ts}>0\) exactly when \(C_{\mathrm{raw},ts}>\vartheta\), increasing \(\vartheta\) can only remove support.

\textbf{Proposition S.3.3:} Let the aggregated column vector be
\(\mathbf m_t=(\mathbf C\mathbf F)_{t,:}^{\top}
=\sum_{s=1}^{T}C_{ts}\mathbf f_s\) and
\begin{equation*}
\LN(\mathbf m_t)
=\boldsymbol{\psi}_{\mathrm{LN}}\odot\hat{\mathbf m}_t
+\mathbf b_{\mathrm{LN}},
\qquad
\|\hat{\mathbf m}_t\|_2\leq\sqrt{d_f}.
\end{equation*}
Here \(\hat{\mathbf m}_t\) is the centered, variance-normalized vector, and \(\boldsymbol\psi_{\mathrm{LN}}\) and \(\mathbf b_{\mathrm{LN}}\) are the learned gain and bias. For finite LayerNorm parameters, define
\begin{equation*}
G_{\mathrm{LN}}
=\|\boldsymbol{\psi}_{\mathrm{LN}}\|_\infty\sqrt{d_f}
+\|\mathbf b_{\mathrm{LN}}\|_2.
\end{equation*}
Then the output in~\eqref{eq:fnsar-output} satisfies
\begin{equation}
\|\mathbf f_t^{\mathrm{out}}\|_2
\leq(1-\zeta)\|\mathbf f_t\|_2+\zeta G_{\mathrm{LN}},
\label{eq:app-norm-bound}
\end{equation}
independent of the scale of \(\mathbf C\).

\emph{Proof:} The LayerNorm bound follows from
\begin{equation*}
\|\LN(\mathbf m_t)\|_2
\leq\|\boldsymbol{\psi}_{\mathrm{LN}}\|_\infty
\|\hat{\mathbf m}_t\|_2+\|\mathbf b_{\mathrm{LN}}\|_2
\leq G_{\mathrm{LN}}.
\end{equation*}
Applying the triangle inequality to~\eqref{eq:fnsar-output} gives
\begin{equation*}
\|\mathbf f_t^{\mathrm{out}}\|_2
\leq(1-\zeta)\|\mathbf f_t\|_2
+\zeta\|\LN(\mathbf m_t)\|_2,
\end{equation*}
and~\eqref{eq:app-norm-bound} follows. This establishes the LayerNorm-residual bound used by the temporal refinement.
